\section{Measurement, constants and output tolerances}
\label{r13-measurement-core}

The existing XOP dimension, frame, sample and dependency types already carry
the information needed here. This chapter adds a physical interpretation and
conditional output certificates, without adding a measurement oracle or a new
packet format. An exactly stored indication, an exact nominal model and the
physical quantity being inferred remain distinct. The purpose is operational:
preserve what is known, keep shared uncertainties coupled, and determine what
an output can claim within a selected application tolerance.

\subsection{Exact defining constants and measured coefficients}
\label{r13-measurement-constants}

Five SI defining constants connect the present time, optical, electrical
and thermal profiles. Their numerical values are fixed by the SI; unit
realizations and measurements made with them still have uncertainty.
The following dimension vectors use the inherited order
$(L,M,T,I,\Theta,N,J_{\rm lum})$.
\href{https://www.bipm.org/en/measurement-units/si-defining-constants}{[ME1: BIPM defining constants]}

\begin{center}\small
\begin{tabularx}{\linewidth}{P{.11\linewidth}P{.45\linewidth}Y}
\toprule
Constant & Exact SI quantity & XOP dimension vector\\
\midrule
$\Delta\nu_{\rm Cs}$ & $9\,192\,631\,770\;\mathrm{Hz}$ & $(0,0,-1,0,0,0,0)$\\
$c$ & $299\,792\,458\;\mathrm{m\,s^{-1}}$ & $(1,0,-1,0,0,0,0)$\\
$h$ & $6.626\,070\,15\times10^{-34}\;\mathrm{J\,s}$ & $(2,1,-1,0,0,0,0)$\\
$e$ & $1.602\,176\,634\times10^{-19}\;\mathrm{C}$ & $(0,0,1,1,0,0,0)$\\
$k_{\rm B}$ & $1.380\,649\times10^{-23}\;\mathrm{J\,K^{-1}}$ & $(2,1,-2,0,-1,0,0)$\\
\bottomrule
\end{tabularx}\end{center}

The decimal coefficients are finite rational literals, reduced canonically by
the XOP grammar. For example, the numerical coefficient of $h$ is
$662607015/10^{42}$ before reduction. A binary floating-point approximation is
unnecessary. The elementary charge $e$ is positive; the charge of an electron
is $-e$. The namespace distinguishes it from the Boolean accepted-event bit.

The gravitational constant $G$ and fine-structure constant $\alpha_{\rm fs}$
are measured quantities, with dimensions $L^3M^{-1}T^{-2}$ and $1$,
respectively. A specialization that needs them binds a dated value, uncertainty
and provenance; this core does not hard-code a new estimate.
\href{https://physics.nist.gov/cuu/Constants/Table/allascii.txt}{[ME2: NIST CODATA table]}
In present SI, vacuum permeability and permittivity are also not independent
exact decimal constants. For positive $\alpha_{\rm fs}$, retain their common
dependency:
\begin{equation}
 \mu_0=\frac{2\alpha_{\rm fs}h}{e^2c},\qquad
 \epsilon_0=\frac{e^2}{2\alpha_{\rm fs}hc},\qquad
 \mu_0\epsilon_0=\frac1{c^2}.
 \label{r13-measurement-vacuum}
\end{equation}
Here $[\mu_0]=MLT^{-2}I^{-2}$ and
$[\epsilon_0]=M^{-1}L^{-3}T^4I^2$. Multiplication cancels the shared
$\alpha_{\rm fs}$ exactly. Replacing the two occurrences by unrelated uncertain
inputs would lose that constraint. The former exact assignment
$\mu_0=4\pi\times10^{-7}\,\mathrm{H\,m^{-1}}$ is not the present SI definition.
\href{https://doi.org/10.59161/AUEZ1291}{[ME3: SI Brochure, 9th edition, 2026, section 2.3.2]}

\subsection{An observation is a modeled acquisition}
\label{r13-measurement-equation}

Let $x(t)$ be the selected physical state, $\theta$ its sensor/calibration
parameters, and $m$ the transduction law into an analogue indication. An
averaging acquisition over $[t_-,t_+]$ has the model
\begin{align}
 M[x,\theta,w]&=\int_{t_-}^{t_+}w(t)\,m(x(t),\theta,t)\,dt,
 &w(t)&\geq0,\quad\int_{t_-}^{t_+}w(t)\,dt=1,\label{r13-measurement-acquisition}\\
 y&=Q_{\Delta}\!\left(M[x,\theta,w]+d_M+\eta\right).
 \label{r13-measurement-observation}
\end{align}
The weighting $w$ has dimension $T^{-1}$; $y$, $M$, the model discrepancy
$d_M$ and analogue noise $\eta$ have the same indication dimension. For a
counter or accumulated dose, use its declared accumulation law instead of a
normalized average. An instantaneous sample is a separately defined profile,
not an implicit substitution of the interval endpoint for the integral.
Each integral requires finite bounds and its stated integrability conditions.
The profile supplies finite function templates or explicit external bindings
for the trajectory and weighting; it cannot encode an arbitrary unknown
continuous signal as a fully known finite literal.

$Q_{\Delta}$ includes its step, offset, ties rule, range and saturation
behavior. For an unsaturated uniform nearest quantizer with step $\Delta>0$,
put $v=M+d_M+\eta$. Then
\begin{equation}
 y=M+d_M+\eta+q,\qquad q=Q_{\Delta}(v)-v,
 \qquad |q|\leq\Delta/2.
 \label{r13-measurement-quantizer}
\end{equation}
This is a deterministic bound, not an assumption that $q$ is uniform,
independent or newly random at every sample. Saturated codes require the
actual preimage of the saturation law; the half-step bound no longer follows.
The stored integer and its rational unit conversion can be exact while
the corresponding physical indication occupies a range.

The measurement model also declares its valid specimen, operating conditions
and reference conditions. A discrepancy bound is supplied evidence or a
stated hypothesis, never a number generated merely by writing $d_M$.
Unrecognized effects are not covered automatically. This makes explicit the
model scope emphasized by JCGM GUM-6, sections 6--7.
\href{https://doi.org/10.59161/JCGMGUM-6-2020}{[ME4: Developing and using measurement models]}

\subsection{Calibration and joint uncertainty}
\label{r13-measurement-calibration}

Let $p$ collect the finite parameters of the chosen profile: initial state,
calibration, clock and frame realization, input history parameters and relevant
residuals. A joint admissible set $\mathcal U$ retains their actual shared
dependencies. With observation map $F$ built from
\eqref{r13-measurement-observation}, the data-compatible set is
\begin{equation}
 \mathcal U_y=\{p\in\mathcal U:F(p)=y\},\qquad
 \mathcal Z_y=\{g(p):p\in\mathcal U_y\}.
 \label{r13-measurement-feasible}
\end{equation}
For multiple records, impose all their observation equations on the same
shared calibration parameters. A repeated common offset does not become a
new independent offset. If the set is proved empty, the declared model,
admissible bounds and observations are inconsistent; empty-set containment
must never be used to certify an output. If solving it is unresolved, retain
the unresolved status. Neither conclusion identifies which physical
assumption failed.

A simple calibration model $y^{\rm cal}=A\theta+\varepsilon$ illustrates what
equations alone can establish. With known design matrix $A\in\R^{m\times r}$,
for consistent noiseless data the parameters are unique exactly when
\begin{equation}
 \operatorname{rank}A=r.
 \label{r13-measurement-rank}
\end{equation}
For a chosen positive-definite covariance $\Sigma$ and full column rank,
the generalized least-squares expression is
\begin{equation}
 \widehat\theta=(A^T\Sigma^{-1}A)^{-1}A^T\Sigma^{-1}y^{\rm cal}.
 \label{r13-measurement-gls}
\end{equation}
The inverse exists since
$z^TA^T\Sigma^{-1}Az=(Az)^T\Sigma^{-1}(Az)>0$ for $z\ne0$.
This proves the algebraic expression's domain, not a calibration accuracy.
Uncertain reference settings make $A$ uncertain and must enter the joint
model. For a differentiable nonlinear calibration, a full-column-rank
Jacobian permits a local regular inverse by selecting independent observation
components. It does not give global uniqueness. Rank deficiency alone does
not decide nonlinear identifiability.

An alternative statistical profile supplies a joint distribution or covariance
with its basis. For a differentiable output map and its linearization
$J=Dg(\widehat p)$,
\begin{equation}
 \Sigma_z\simeq J\Sigma_pJ^T,\qquad
 u^2(z_i)\simeq\sum_{j,k}J_{ij}(\Sigma_p)_{jk}J_{ik}.
 \label{r13-measurement-covariance}
\end{equation}
Cross terms retain correlations. This is exact for an affine map with the
stated covariance; in general it is a first-order approximation and a
standard uncertainty is not a worst-case bound.
\href{https://doi.org/10.59161/JCGM100-2008E}{[ME5: JCGM 100, sections 5.1--5.2]}
Significant nonlinearity can also affect the output estimate itself;
exact evaluation of $g(\widehat p)$ does not make it the distributional mean.
\href{https://doi.org/10.59161/PPDI3267}{[ME6: JCGM 100, Amendment 1:2026]}

\subsection{Acquisition time and realized coordinates}
\label{r13-measurement-time}

R12's WI record, immutable registry and session consumption rules remain
unchanged. The calibration identifier resolves complete calibration content;
its presence is not proof of freshness or traceability. Add the chosen
measurement model, acquisition weighting, quantizer and joint-uncertainty
definition to that existing dependency closure.

For recorded source times $\tau_\pm$, a declared affine clock model is
\begin{equation}
 t_\pm=a\tau_\pm+b+\delta t_\pm,\qquad a>0.
 \label{r13-measurement-clock}
\end{equation}
The same $a,b$ appear at both endpoints and in every record covered by that
calibration. The physical acquisition interval, conversion validity interval
and availability time remain distinct. Subtracting endpoints cancels the
shared offset $b$; treating it as two independent offsets would miss this
exact dependency. The uncertainty of $a$ and endpoint residuals remains.

The analogous physical frame map is
\begin{equation}
 x_B(t)=R_{BA}(t)x_A(t)+o_{BA}(t),\qquad
 R_{BA}^TR_{BA}=I,\quad\det R_{BA}=1.
 \label{r13-measurement-frame}
\end{equation}
The XOP frame identity specifies the convention; $R_{BA}$ and $o_{BA}$
specify its realization. A shared uncertain mounting is retained once in
$p$, and an interval enclosure of its entries must retain, or conservatively
enclose, the rotation constraints. A nominal matrix is not made physically
aligned merely by encoding it exactly. Coordinate values from different
acquisition times require the selected dynamical/time-transfer model before
they can represent simultaneous geometry.

\subsection{A bounded output certificate}
\label{r13-measurement-certificate}

The core needs a reusable rule for deciding a functional tolerance without
inventing a universal numerical tolerance. Let $K$ be a nonempty compact
convex set in the finite parameter coordinates; require
$\widehat p\in K$ and a nonempty $\mathcal U_y\subseteq K$. The output
$g_i$ is continuously differentiable on an open neighborhood of $K$. Suppose
the following finite bounds have been established for all $p\in K$:
\begin{equation}
 |p_j-\widehat p_j|\leq r_j,\qquad
 \left|\frac{\partial g_i}{\partial p_j}(p)\right|\leq L_{ij}.
 \label{r13-measurement-lipschitz-data}
\end{equation}
$L_{ij}$ has output-unit per parameter-unit dimensions. No covariance or
independence assumption is used. By integrating the derivative along the
segment from $\widehat p$ to $p$,
\begin{align}
 g_i(p)-g_i(\widehat p)
 &=\int_0^1\sum_j
   \frac{\partial g_i}{\partial p_j}
   (\widehat p+s(p-\widehat p))(p_j-\widehat p_j)\,ds,\\
 |g_i(p)-g_i(\widehat p)|&\leq\sum_jL_{ij}r_j.
 \label{r13-measurement-meanvalue}
\end{align}
The compact domain stays clear of undefined divisions, branch changes or
other points where the derivative hypotheses fail. A box can conservatively
enclose a correlated set; using the original joint constraints can tighten
the result. Piecewise or discrete outputs use direct image/branch enclosure
instead of differentiating through a switch.

To apply the certificate to a physical result, assume its actual parameter
vector lies in $\mathcal U_y$. Nonemptiness alone does not establish that
physical inclusion. Suppose the intended physical output is
$z_i^{\rm phys}=g_i(p)+d_i^{\rm out}$ with
$|d_i^{\rm out}|\leq b_i$, and the delivered numerical output obeys
$|z_i^{\rm del}-g_i(\widehat p)|\leq\rho_i$. Then
\begin{equation}
 |z_i^{\rm phys}-z_i^{\rm del}|\leq
 B_i:=\sum_jL_{ij}r_j+b_i+\rho_i.
 \label{r13-measurement-total-bound}
\end{equation}
This triangle-inequality certificate separates input/calibration uncertainty,
output-model discrepancy and representation/evaluation error. Effects already
included in $p$ are not counted again in $b_i$. The XOP exact expression can
be retained with no numerical rendering error in its denotation; a finite
rendering needs its own $\rho_i$ certificate. Existing binary64 execution is
not assigned $\rho_i=0$ by this new formalization.

For a requested absolute error tolerance $\tau_i\geq0$, the condition
$B_i\leq\tau_i$ is sufficient on the stated domain and acquisition interval.
For an allowed physical operating band $[\ell_i,u_i]$, sufficient containment is
\begin{equation}
 z_i^{\rm del}-B_i\geq\ell_i,\qquad
 z_i^{\rm del}+B_i\leq u_i.
 \label{r13-measurement-tolerance}
\end{equation}
Failure to prove this containment does not prove violation: a conservative
enclosure may overlap both regions. A wholly disjoint enclosing interval
proves exclusion; overlap remains unresolved. This is the general form of
R12's bounded bridge and timing decisions. It certifies the declared quantity
and interval, not every later state or a different physical application.

\subsection{Minimal specialization record}
\label{r13-measurement-record}

Each specialization binds only: the measurand and conditions; the complete
measurement/calibration expressions and actual source records; the unit,
frame and acquisition-time realizations; the joint uncertainty or statistical
model; and the requested output tolerance with its scoped evidence. The
certificate binds those same dependency versions. Changing a calibration,
clock, model basis, duration or tolerance requires reevaluating the affected
claim, not changing the one-bit packing format.

These are metadata and template/dependency obligations under XOP-R1. Set
images, universal bounds and physical evidence are not secretly new executable
opcodes. A backend may return \code{UNKNOWN}, \code{MISSING\_INPUT} or
\code{RESOURCE\_LIMIT} when it cannot establish an obligation. The bridge,
thermal and optical expressions remain useful before that backend exists.

Metrological traceability is a property of a measurement result supported by
its calibration chain and uncertainty; source hashing and exact arithmetic
preserve provenance but do not supply that chain. Fitness for a particular
use additionally depends on the required uncertainty or tolerance.
\href{https://www.nist.gov/calibrations/traceability}{[ME7: NIST traceability policy]}
The additions here make that distinction actionable through
\eqref{r13-measurement-total-bound}, rather than adding unrelated domain rules.
